% GENERATED FILE. Edit canonical lessons, then run npm run generate:books.
% canonical-source-hash: fnv1a64:147b8efeeda6f56c
% canonical-lessons: SA-C482-cikitsati, SA-C482-aturah, SA-C482-nidanam, SA-C482-upacarah, SA-C482-salyam

\chapter{At the Physician's}
\label{ch:sa-at-the-physician-s}
\hlchaptermodality{482}

\begin{chapteropening}
\textbf{By the end of this chapter:} \emph{I can talk about a visit to the physician.}

Complete the last lesson of chapter 482: I can talk about a visit to the physician.
\end{chapteropening}

\section[\texorpdfstring{\sk{चिकित्सति} (cikitsati)}{चिकित्सति (cikitsati)}]{\sk{चिकित्सति} (cikitsati) — treats a patient}

\label{lesson:SA-C482-cikitsati}



\begin{quote}
Before the new one: say the Sanskrit for lame, then the Sanskrit for hunchbacked.
\end{quote}

\subsection*{You'll want to know: \sk{चिकित्सति}}
\textbf{\sk{चिकित्सति}} — \emph{cikitsati} — \textquotedblleft{}treats a patient\textquotedblright{}.

\begin{quote}
\sk{वैद्यः} \sk{रुग्णं} \sk{बालं} \sk{चिकित्सति।}
\end{quote}

\begin{quote}
\emph{vaidyaḥ rugṇaṃ bālaṃ cikitsati.}
\end{quote}

\begin{quote}
The physician treats the sick child.
\end{quote}

Past (he …): \sk{सः} \sk{रुग्णं} \sk{बालं} \sk{चिकित्सितवान्}, \emph{saḥ rugṇaṃ bālaṃ cikitsitavān}.

\begin{cousinweb}
\textbf{\sk{चिकित्सति}} (\emph{cikitsati}) is built on the desiderative of the root √cit, \textquotedblleft{}to perceive, to care for\textquotedblright{}, the word behind \textbf{\sk{चिकित्सकः}}; MW gives \textquotedblleft{}to treat medically, cure\textquotedblright{} (MW).
\end{cousinweb}

\begin{grammarlens}[title={What you've built}]
One more word for describing what you lived and felt.
\end{grammarlens}

\subsection*{Your turn}
\begin{itemize}
\raggedright
  \item \emph{Say it:} \emph{cikitsati}
  \item \emph{Say it:} \emph{cikitsati}, once more
  \item \emph{Say it:} \emph{vaidyaḥ rugṇaṃ bālaṃ cikitsati.} — The physician treats the sick child.
  \item \emph{Recall:} say the Sanskrit for lame, then the Sanskrit for hunchbacked, then say \emph{cikitsati} again
\end{itemize}

\begin{tcolorbox}[breakable,colback=teal!4,colframe=teal!35!black,title={Before you move on}]
What does \sk{चिकित्सति} mean? (\textquotedblleft{}Treats a patient\textquotedblright{}.) Say it once more.
\end{tcolorbox}
\section[\texorpdfstring{\sk{आतुरः} (āturaḥ)}{आतुरः (āturaḥ)}]{\sk{आतुरः} (āturaḥ) — a patient; sick}

\label{lesson:SA-C482-aturah}



\begin{quote}
Before the new one: say the Sanskrit for hunchbacked, then the Sanskrit for treats a patient.
\end{quote}

\subsection*{You'll want to know: \sk{आतुरः}}
\textbf{\sk{आतुरः}} — \emph{āturaḥ} — \textquotedblleft{}a patient; sick\textquotedblright{}.

\begin{quote}
\sk{वैद्यः} \sk{आतुरस्य} \sk{नाडीं} \sk{स्पृशति।}
\end{quote}

\begin{quote}
\emph{vaidyaḥ āturasya nāḍīṃ spṛśati.}
\end{quote}

\begin{quote}
The physician feels the patient's pulse.
\end{quote}

A masculine or feminine noun, by the person meant.

\begin{cousinweb}
\textbf{\sk{आतुरः}} (\emph{āturaḥ}) MW divides it \emph{ā-tura} and gives \textquotedblleft{}suffering, sick (in body or mind)\textquotedblright{}; it names no root, so its origin is uncertain (MW).
\end{cousinweb}

\begin{grammarlens}[title={What you've built}]
One more word for describing what you lived and felt.
\end{grammarlens}

\subsection*{Your turn}
\begin{itemize}
\raggedright
  \item \emph{Say it:} \emph{āturaḥ}
  \item \emph{Say it:} \emph{āturaḥ}, once more
  \item \emph{Say it:} \emph{vaidyaḥ āturasya nāḍīṃ spṛśati.} — The physician feels the patient's pulse.
  \item \emph{Recall:} say the Sanskrit for hunchbacked, then the Sanskrit for treats a patient, then say \emph{āturaḥ} again
\end{itemize}

\begin{tcolorbox}[breakable,colback=teal!4,colframe=teal!35!black,title={Before you move on}]
What does \sk{आतुरः} mean? (\textquotedblleft{}A patient; sick\textquotedblright{}.) Say it once more.
\end{tcolorbox}
\section[\texorpdfstring{\sk{निदानम्} (nidānam)}{निदानम् (nidānam)}]{\sk{निदानम्} (nidānam) — the cause of an illness}

\label{lesson:SA-C482-nidanam}



\begin{quote}
Before the new one: say the Sanskrit for treats a patient, then the Sanskrit for a patient.
\end{quote}

\subsection*{You'll want to know: \sk{निदानम्}}
\textbf{\sk{निदानम्}} — \emph{nidānam} — \textquotedblleft{}the cause of an illness\textquotedblright{}.

\begin{quote}
\sk{वैद्यः} \sk{रोगस्य} \sk{निदानं} \sk{जानाति।}
\end{quote}

\begin{quote}
\emph{vaidyaḥ rogasya nidānaṃ jānāti.}
\end{quote}

\begin{quote}
The physician knows the cause of the illness.
\end{quote}

A neuter noun.

\begin{cousinweb}
\textbf{\sk{निदानम्}} (\emph{nidānam}) is built from \emph{ni} and the root √dā, \textquotedblleft{}to bind\textquotedblright{}: MW's first sense is \textquotedblleft{}a band, a halter\textquotedblright{}, then \textquotedblleft{}a first cause\textquotedblright{} and \textquotedblleft{}the cause of a disease and enquiry into it, pathology\textquotedblright{} (MW).
\end{cousinweb}

\begin{grammarlens}[title={What you've built}]
One more word for describing what you lived and felt.
\end{grammarlens}

\subsection*{Your turn}
\begin{itemize}
\raggedright
  \item \emph{Say it:} \emph{nidānam}
  \item \emph{Say it:} \emph{nidānam}, once more
  \item \emph{Say it:} \emph{vaidyaḥ rogasya nidānaṃ jānāti.} — The physician knows the cause of the illness.
  \item \emph{Recall:} say the Sanskrit for treats a patient, then the Sanskrit for a patient, then say \emph{nidānam} again
\end{itemize}

\begin{tcolorbox}[breakable,colback=teal!4,colframe=teal!35!black,title={Before you move on}]
What does \sk{निदानम्} mean? (\textquotedblleft{}The cause of an illness\textquotedblright{}.) Say it once more.
\end{tcolorbox}
\section[\texorpdfstring{\sk{उपचारः} (upacāraḥ)}{उपचारः (upacāraḥ)}]{\sk{उपचारः} (upacāraḥ) — medical treatment}

\label{lesson:SA-C482-upacarah}



\begin{quote}
Before the new one: say the Sanskrit for a patient, then the Sanskrit for the cause of an illness.
\end{quote}

\subsection*{You'll want to know: \sk{उपचारः}}
\textbf{\sk{उपचारः}} — \emph{upacāraḥ} — \textquotedblleft{}medical treatment\textquotedblright{}.

\begin{quote}
\sk{उचितेन} \sk{उपचारेण} \sk{आतुरः} \sk{स्वस्थः} \sk{भवति।}
\end{quote}

\begin{quote}
\emph{ucitena upacāreṇa āturaḥ svasthaḥ bhavati.}
\end{quote}

\begin{quote}
With the right treatment the patient gets well.
\end{quote}

A masculine noun.

\begin{cousinweb}
\textbf{\sk{उपचारः}} (\emph{upacāraḥ}) is built from \emph{upa}, \textquotedblleft{}near\textquotedblright{}, and the root √car, \textquotedblleft{}to move\textquotedblright{}; MW gives \textquotedblleft{}attendance, service\textquotedblright{} and \textquotedblleft{}attendance on a patient, medical practice\textquotedblright{} (MW).
\end{cousinweb}

\begin{grammarlens}[title={What you've built}]
One more word for describing what you lived and felt.
\end{grammarlens}

\subsection*{Your turn}
\begin{itemize}
\raggedright
  \item \emph{Say it:} \emph{upacāraḥ}
  \item \emph{Say it:} \emph{upacāraḥ}, once more
  \item \emph{Say it:} \emph{ucitena upacāreṇa āturaḥ svasthaḥ bhavati.} — With the right treatment the patient gets well.
  \item \emph{Recall:} say the Sanskrit for a patient, then the Sanskrit for the cause of an illness, then say \emph{upacāraḥ} again
\end{itemize}

\begin{tcolorbox}[breakable,colback=teal!4,colframe=teal!35!black,title={Before you move on}]
What does \sk{उपचारः} mean? (\textquotedblleft{}Medical treatment\textquotedblright{}.) Say it once more.
\end{tcolorbox}
\section[\texorpdfstring{\sk{शल्यम्} (śalyam)}{शल्यम् (śalyam)}]{\sk{शल्यम्} (śalyam) — a splinter}

\label{lesson:SA-C482-salyam}



\begin{quote}
Before the new one: say the Sanskrit for the cause of an illness, then the Sanskrit for medical treatment.
\end{quote}

\subsection*{You'll want to know: \sk{शल्यम्}}
\textbf{\sk{शल्यम्}} — \emph{śalyam} — \textquotedblleft{}a splinter\textquotedblright{}.

\begin{quote}
\sk{पादे} \sk{शल्यम्} \sk{अस्ति}, \sk{अतः} \sk{सः} \sk{शनैः} \sk{चलति।}
\end{quote}

\begin{quote}
\emph{pāde śalyam asti, ataḥ saḥ śanaiḥ calati.}
\end{quote}

\begin{quote}
There is a splinter in his foot, so he walks slowly.
\end{quote}

A neuter noun.

\begin{cousinweb}
\textbf{\sk{शल्यम्}} (\emph{śalyam}) MW's first sense is \textquotedblleft{}a dart, a lance\textquotedblright{}, then \textquotedblleft{}anything causing pain\textquotedblright{} and, in medicine, \textquotedblleft{}anything lodged in the body, a splinter\textquotedblright{}; MW gives no derivation, so its origin is uncertain (MW).
\end{cousinweb}

\begin{grammarlens}[title={What you've built}]
One more word for describing what you lived and felt.
\end{grammarlens}

\subsection*{Your turn}
\begin{itemize}
\raggedright
  \item \emph{Say it:} \emph{śalyam}
  \item \emph{Say it:} \emph{śalyam}, once more
  \item \emph{Say it:} \emph{pāde śalyam asti, ataḥ saḥ śanaiḥ calati.} — There is a splinter in his foot, so he walks slowly.
  \item \emph{Recall:} say the Sanskrit for the cause of an illness, then the Sanskrit for medical treatment, then say \emph{śalyam} again
\end{itemize}

\begin{tcolorbox}[breakable,colback=teal!4,colframe=teal!35!black,title={Before you move on}]
What does \sk{शल्यम्} mean? (\textquotedblleft{}A splinter\textquotedblright{}.) Say it once more.
\end{tcolorbox}
